\section{Introduction}
\label{sec:intro}

For a developer who needs to run a small HTTP service on AWS, the
choice of compute primitive is uncomfortable. EC2~\cite{aws-ec2} gives
predictable performance and the lowest unit price once an instance is
busy, but shifts capacity planning and OS upkeep onto the user.  ECS
Fargate~\cite{aws-fargate} hides that work but bills per vCPU-second of
running task time. AWS Lambda~\cite{aws-lambda} bills per request and
per GB-second with no idle cost, but introduces a cold-start penalty
and a per-invocation overhead absent on long-running workers.

Existing studies illuminate one part of this picture each.
Wang et al.~\cite{peeking-serverless} characterise FaaS runtime behaviour,
McGrath and Brenner~\cite{mcgrath-faas} compare cold-start strategies
across early serverless platforms, and Eismann
et al.~\cite{eismann2021review} review serverless performance studies.
None of these compare the \emph{three native AWS deployment
families} side-by-side under a workload that is identical down to the
container image. That is what we set out to do.

\paragraph{Methodology in one sentence.} We package a stateless URL
shortener (\texttt{FastAPI} + DynamoDB) as a single OCI container image,
attach the AWS Lambda Web Adapter so the same image runs on a Lambda
container function as well as on EC2 and Fargate, and put all three
deployments behind a shared Application Load Balancer. A
\texttt{t3.small} k6 worker in the same availability zone exercises the
endpoints under four scenarios.

\paragraph{Contributions.} The work makes four contributions:
\begin{enumerate}[leftmargin=*,topsep=2pt,itemsep=2pt]
  \item A reproducible benchmarking harness in which the same image is
        deployed to EC2, Fargate, and Lambda from the same Terraform
        codebase, runnable end-to-end on any AWS account
        environment that the course uses (\cref{sec:design}).
  \item A measured comparison of throughput, p50/p95/p99 latency, and
        elasticity across the three paradigms with bootstrap confidence
        intervals on every percentile (\cref{sec:results-steady}).
  \item A characterisation of Lambda cold-start behaviour for a 1\,GB
        container function and its impact on burst latency
        (Sections~\ref{sec:results-elasticity} and~\ref{sec:results-cold}).
  \item An analytic cost model, calibrated against the measured
        per-request latency, that yields explicit break-even points
        between the three paradigms and a Pareto front of
        cost-versus-latency choices (\cref{sec:results-cost}).
\end{enumerate}

\paragraph{Headline findings.} (i)~At every load level we sustained,
EC2 keeps the tightest tail latency: its p99 sits under
60\,ms at 500\,RPS while Lambda's exceeds 200\,ms even after warm-up.
(ii)~Fargate scales out smoothly under a step-arrival burst and reaches
steady-state in roughly 25\,s, against $\sim$45\,s for EC2 and
$<$10\,s for Lambda once a few additional concurrent containers are
spawned. (iii)~Lambda's first-request latency after a cold container
spawn is dominated by image initialisation; we measured a median
of $\sim$1.4\,s and a p99 of $\sim$11\,s for our 1\,GB container.
(iv)~The cost cross-over points implied by our calibrated model put
Lambda ahead of Fargate up to $\sim$17\,RPS and ahead of EC2 up to
$\sim$60\,RPS for a workload of our latency profile.

The remainder of the paper is organised as follows.
\Cref{sec:background} situates the three paradigms and their billing.
\Cref{sec:design} describes the application and the unified deployment
pipeline. \Cref{sec:methodology} explains the workload, metrics, and
statistical procedure. \Cref{sec:results} reports steady-state,
elastic, cold-start, and cost results.
Sections~\ref{sec:discussion}--\ref{sec:conclusion} discuss
limitations, related work, and conclude.
